\chapter{Metric comparison and uniform Weyl asymptotics}\label{sec:width-metrics}
Throughout, $M^{n+1}$ is smooth, closed and connected, $n\geq1$, and $\omega_p$ is the mod-two volume spectrum of Definition \ref{def:p-width}. We compare smooth Riemannian metrics on the same manifold. The normalized metric estimate follows the argument in \cite{MNS19}, Lemma 1, p.~425. The Weyl theorem itself, Theorem \ref{thm:weyl-volume}, remains external.

\begin{lemma}[Mass and flat norms under metric comparison]\label{lem:metric-mass-flat-comparison}
If $a^2g\leq h\leq b^2g$ as quadratic forms, where $0<a\leq b$, then every rectifiable mod-two $k$-chain $T$ satisfies
\[
 a^k\mathbf M_g(T)\leq\mathbf M_h(T)\leq b^k\mathbf M_g(T).
\]
For the flat norm of an $n$-chain,
\[
 \min\{a^n,a^{n+1}\}\mathcal F_g(T)
 \leq\mathcal F_h(T)
 \leq\max\{b^n,b^{n+1}\}\mathcal F_g(T).
\]
In particular the two flat topologies agree.
\end{lemma}
\begin{proof}
On an approximate tangent $k$-plane the eigenvalues of $h$ relative to $g$ lie in $[a^2,b^2]$. The square root of their product, the area-density ratio, lies in $[a^k,b^k]$. Integrate against the rectifiable multiplicity to obtain the mass bounds, also on every measurable subset. For each flat decomposition $T=R+\partial S$, with $R$ of dimension $n$ and $S$ of dimension $n+1$, apply the upper mass bounds to the two terms and take the infimum. This gives the upper flat bound. Interchanging the two metrics gives the lower bound. Smooth uniformly equivalent metrics have the same rectifiable and finite-mass chain classes, so these infima use the same decompositions.
\end{proof}

\begin{proposition}[Sweepout admissibility is independent of a smooth metric]\label{prop:metric-sweepout-admissibility}
An admissible $p$-sweepout for one smooth metric on $M$ is admissible for every other smooth metric on $M$, with the same detecting cohomology class.
\end{proposition}
\begin{proof}
Compactness of $M$ gives comparison constants $a,b$ as in Lemma \ref{lem:metric-mass-flat-comparison}. The same underlying cycle group has the same flat topology, boundary component and distinguished nonzero degree-one class under the identity map. Flat continuity and cohomological detection are therefore unchanged. Maximal mass remains bounded by the mass inequality. Distances satisfy $d_h\geq a d_g$, so $B_r^h(y)\subset B_{r/a}^g(y)$, and the localized mass inequality gives
\[
 \|\Phi(x)\|_h(B_r^h(y))
 \leq b^n\|\Phi(x)\|_g(B_{r/a}^g(y)).
\]
Taking the supremum and then $r\downarrow0$ proves no concentration for $h$. The reverse comparison proves equivalence.
\end{proof}

\begin{proposition}[Bilateral metric bounds for every width]\label{prop:width-metric-comparison}
If $a^2g\leq h\leq b^2g$, then for every $p\geq1$,
\[
 a^n\omega_p(M,g)\leq\omega_p(M,h)\leq b^n\omega_p(M,g).
\]
This is the metric comparison used in \cite{MNS19}, proof of Lemma 1.
\end{proposition}
\begin{proof}
Proposition \ref{prop:metric-sweepout-admissibility} identifies the two sets of competitors. Apply the $n$-mass comparison to every slice of each family, take the supremum in the parameter, and then take the infimum over families. This proves both bounds. In particular $g\leq h$ implies $\omega_p(M,g)\leq\omega_p(M,h)$. Taking $a=b$ recovers the scaling formula of Lemma \ref{lem:width-scaling}.
\end{proof}

\begin{proposition}[Diffeomorphism invariance of the spectrum]\label{prop:width-diffeomorphism-invariance}
If $F:M\to M$ is a smooth diffeomorphism and $g$ is smooth, then
\[
 \omega_p(M,F^*g)=\omega_p(M,g)\quad(p\geq1).
\]
No orientation assumption on $F$ is required for mod-two coefficients.
\end{proposition}
\begin{proof}
The map $F:(M,F^*g)\to(M,g)$ is an isometry. Its pushforward on mod-two chains preserves boundary, masses, flat norms, and metric-ball mass estimates. Thus it gives a bijection of the admissible families. Its induced homeomorphism of the boundary-cycle component preserves the unique nonzero degree-one generator and hence all of its cup powers. Corresponding families have identical maximal masses, so their infima agree. The pushforward uses mod-two multiplicities and therefore does not require a choice of ambient orientation.
\end{proof}

\begin{theorem}[Uniform Lipschitz control of normalized widths]\label{thm:normalized-width-lipschitz}
Let $M^{n+1}$ be smooth, closed and connected, with $n\geq1$. Fix a smooth metric $g_0$ on $M$ and $0<A\leq B$. There is $K=K(M,g_0,A,B,n)$ such that all smooth metrics $g,h$ with $Ag_0\leq g,h\leq Bg_0$ satisfy
\[
 \left|p^{-1/(n+1)}\omega_p(M,g)-p^{-1/(n+1)}\omega_p(M,h)\right|
 \leq K\|g-h\|_{C^0(g_0)}
\]
for every $p\geq1$. Here the norm is $\sup_{x,v:\,g_0(v,v)=1}|(g-h)(v,v)|$. This is the smooth-metric form of \cite{MNS19}, Lemma 1.
\end{theorem}
\begin{proof}
Write $\varepsilon=\|g-h\|_{C^0(g_0)}$. Since $g\geq Ag_0$, we have $h\leq(1+\varepsilon/A)g$. Proposition \ref{prop:width-metric-comparison} gives
\[
 \omega_p(h)-\omega_p(g)
 \leq\left((1+\varepsilon/A)^{n/2}-1\right)\omega_p(g)
 \leq\left((1+\varepsilon/A)^{n/2}-1\right)B^{n/2}\omega_p(g_0).
\]
The norm is at most $2B$ on this band. The derivative of $t\mapsto(1+t)^{n/2}$ is bounded on $[0,2B/A]$, so the parenthesis is at most a constant times $\varepsilon$. Reverse $g,h$ for the other inequality. Finally Theorem \ref{thm:weyl-volume} bounds $p^{-1/(n+1)}\omega_p(g_0)$ for large $p$. The finitely many smaller widths are finite and bounded by a sufficiently large one using Lemma \ref{lem:width-monotonicity}. This gives a bound uniform in all $p$. This last step uses the external Weyl theorem; the source instead invokes the external Gromov--Guth growth bound, \cite{LMN}, Theorem 2.8.
\end{proof}

\begin{theorem}[The Weyl limit is uniform on compact metric families]\label{thm:compact-metric-uniform-weyl}
Let $\mathcal G$ be a $C^0$-compact set of smooth Riemannian metrics on the fixed closed connected $M^{n+1}$, $n\geq1$. Then
\[
 \lim_{p\to\infty}\sup_{g\in\mathcal G}
 \left|p^{-1/(n+1)}\omega_p(M,g)
 -a(n)\operatorname{Vol}(M,g)^{n/(n+1)}\right|=0.
\]
This is a local uniformity consequence of the pointwise Weyl law in \cite{LMN}, Theorem 1.1, and the estimate of \cite{MNS19}, Lemma 1.
\end{theorem}
\begin{proof}
Fix $g_0$. Compactness and positive definiteness give common constants $A,B>0$ with $Ag_0\leq g\leq Bg_0$ for every $g\in\mathcal G$. Indeed the smallest and largest relative eigenvalues vary continuously with the metric in $C^0$, and their extrema over $M$ and the compact family are positive and finite. The functions $W_p(g)=p^{-1/(n+1)}\omega_p(M,g)$ have a common Lipschitz constant by Theorem \ref{thm:normalized-width-lipschitz}. The candidate limit $W(g)=a(n)\operatorname{Vol}(M,g)^{n/(n+1)}$ is continuous: the volume density is the square root of the metric determinant, uniformly bounded above and below on this band.
Given $\varepsilon>0$, choose a finite sufficiently fine $C^0$ net $g_1,\ldots,g_N$ in $\mathcal G$ so both the common Lipschitz error and the uniform-continuity error for $W$ are below $\varepsilon/3$. The pointwise external Weyl law makes $|W_p(g_i)-W(g_i)|<\varepsilon/3$ for all net points once $p$ is large. The triangle inequality then bounds $|W_p(g)-W(g)|$ by $\varepsilon$ for every $g\in\mathcal G$. No claim about nonsmooth limiting metrics or derivatives is made.
\end{proof}

\begin{proposition}[Volume normalization removes constant scaling]\label{prop:volume-normalized-width}
For each $p\geq1$, the quantity
\[
 \mathcal W_p(g)=\frac{\omega_p(M,g)}{\operatorname{Vol}(M,g)^{n/(n+1)}}
\]
is unchanged by constant metric scaling $g\mapsto r^2g$ and by pullback under a diffeomorphism. It is continuous in the $C^0$ topology on smooth metrics. Moreover $p^{-1/(n+1)}\mathcal W_p(g)\to a(n)$ uniformly on every $C^0$-compact family of smooth metrics.
\end{proposition}
\begin{proof}
Lemma \ref{lem:width-scaling} multiplies the numerator and denominator by the same factor $r^n$. Diffeomorphisms preserve volume under pullback, and Proposition \ref{prop:width-diffeomorphism-invariance} handles the numerator. Theorem \ref{thm:normalized-width-lipschitz} gives continuity of the numerator for fixed $p$ on each uniform metric band; the determinant formula gives continuity and positivity of the denominator. On a compact family that denominator is bounded away from zero. Divide the uniform estimate in Theorem \ref{thm:compact-metric-uniform-weyl} by it to obtain the final assertion. This normalized quantity still depends on the shape of the metric; scaling invariance is not a statement that volume determines a finite width.
\end{proof}

\begin{proposition}[First variation with respect to the ambient metric]\label{prop:metric-area-volume-variation}
Let $g_t$ be a smooth family of metrics with $g_0=g$ and $\dot g_0=h$. For a fixed smooth closed immersed $n$-dimensional hypersurface $\Sigma$,
\[
 \left.\frac d{dt}\right|_0\operatorname{Area}_{g_t}(\Sigma)
 =\frac12\int_\Sigma\operatorname{tr}_{T\Sigma,g}h\,dA_g,
 \qquad
 \left.\frac d{dt}\right|_0\operatorname{Vol}(M,g_t)
 =\frac12\int_M\operatorname{tr}_g h\,dV_g.
\]
For $g_t=e^{2tf}g$, these become $n\int_\Sigma f\,dA_g$ and $(n+1)\int_M f\,dV_g$. The hypersurface formula is the metric derivative appearing in \cite{MNS19}, Lemma 2, p.~426.
\end{proposition}
\begin{proof}
In fixed local coordinates on $\Sigma$, let $G(t)$ be the induced metric matrix. The determinant identity gives
\[
 \frac d{dt}\sqrt{\det G(t)}\Big|_0
 =\frac12\sqrt{\det G(0)}\operatorname{tr}(G(0)^{-1}\dot G(0)).
\]
The trace is precisely the restriction of $h$ to the tangent space, traced with its induced metric. A partition of unity and compactness permit differentiation under the integral. The same calculation in ambient coordinates proves the volume formula. For the conformal variation, $h=2fg$ has tangential trace $2nf$ and ambient trace $2(n+1)f$. The calculation is local and does not need orientability of the hypersurface.
\end{proof}

\begin{proposition}[Width differentiation along an assumed smooth realizing family]\label{prop:smooth-width-realizer-derivative}
Suppose $g_t$ is smooth for $t$ near zero and there are finitely many smoothly varying closed embedded $g_t$-minimal hypersurfaces $\Sigma_i(t)$ and fixed positive integers $m_i$ such that
\[
 \omega_p(M,g_t)=\sum_i m_i\operatorname{Area}_{g_t}(\Sigma_i(t))
\]
throughout that interval. Then the width is differentiable at zero and
\[
 \left.\frac d{dt}\right|_0\omega_p(M,g_t)
 =\frac12\sum_i m_i\int_{\Sigma_i(0)}\operatorname{tr}_{T\Sigma_i(0),g_0}\dot g_0\,dA_{g_0}.
\]
The existence of such a smooth realizing family is an assumption here. The generic-family derivative representation in \cite{MNS19}, Lemma 2, is a separate theorem with additional regularity, transversality and dimension hypotheses.
\end{proposition}
\begin{proof}
Differentiate each area as a function of both the metric and the immersion. Its metric contribution is Proposition \ref{prop:metric-area-volume-variation}. Its immersion contribution at fixed $g_0$ is the first variation of the closed minimal hypersurface and vanishes, for every variation field, whether or not the normal bundle is trivial. The finite weighted sum is therefore differentiable with the displayed derivative. By the assumed identity on an entire interval this derivative equals that of the width. An equality at a single time would not justify this step, and pointwise existence of a minimizing varifold would not give a smooth family.
\end{proof}

\begin{example}[Uniform Weyl convergence does not justify differentiating its limit]\label{ex:uniform-limit-derivatives}
Uniform convergence, even with a common Lipschitz bound, does not imply convergence of derivatives at a point where all the functions and their limit are differentiable. Consequently Theorem \ref{thm:compact-metric-uniform-weyl} alone cannot supply a derivative formula for widths or equidistribution of individual realizing hypersurfaces.
\end{example}
\begin{proof}
On $[-1,1]$ take $f_j(t)=\sin(jt)/j$. Then $\|f_j\|_\infty\leq1/j$, so $f_j\to0$ uniformly, and $|f_j'(t)|=|\cos(jt)|\leq1$. Nevertheless $f_j'(0)=1$ for every $j$, whereas the derivative of the limit is zero. This is an elementary logical obstruction to exchanging a limit and differentiation; it does not assert that these particular functions are width functions. The additional geometric and nonsmooth variational arguments in \cite{MNS19}, Lemmas 2--3, are not replaced by uniform convergence.
\end{proof}
